\documentclass[sigplan]{acmart}

\usepackage{tikz}

\setcopyright{none}

\begin{document}

\title{How to Add the ACM Badges to Your Paper}

\author{Jungwon Kim}
\orcid{0000-0001-6594-6225}
\affiliation{\institution{Oak Ridge National Laboratory}\city{Oak Ridge}\state{TN}\country{USA}}
\email{kimj@ieee.org}

\begin{abstract}
ACM keywords are mandatory for papers over two pages.
CCS concepts are mandatory for papers over two pages.
\end{abstract}

\begin{CCSXML}
<ccs2012>
<concept>
<concept_id>10011007.10011006.10011008</concept_id>
<concept_desc>Software and its engineering~General programming languages</concept_desc>
<concept_significance>500</concept_significance>
</concept>
</ccs2012>
\end{CCSXML}

\ccsdesc[500]{Software and its engineering~General programming languages}

\keywords{High-performance computing}

\maketitle

\begin{tikzpicture}[remember picture, overlay]
\node[anchor=north east, inner xsep=50pt, inner ysep=10pt] at (current page.north east) {
\href{https://www.acm.org/publications/policies/artifact-review-and-badging-current}{
\includegraphics[width=50pt]{artifacts_evaluated_functional_v1.1.pdf}
\includegraphics[width=50pt]{artifacts_evaluated_reusable_v1.1.pdf}
\includegraphics[width=50pt]{artifacts_available_v1.1.pdf}
\includegraphics[width=50pt]{results_reproduced_v1.1.pdf}
\includegraphics[width=50pt]{results_replicated_v1.1.pdf}
}};
\end{tikzpicture}

\section{Instructions}
The badges are placed on the top right corner of the paper, side-by-side. The “most important” badge is positioned rightmost, with the second most important to the left of it, and so on. The badges are listed from least to most important on the badging page~\cite{ACM} (so, pink/Functional is least important, light blue/Replicated is most important).

\bibliographystyle{ACM-Reference-Format}
\bibliography{sample}

\end{document}

